\documentclass[11pt]{article}
\usepackage{amsmath,amssymb,amsthm}
\usepackage[margin=1.1in]{geometry}

\newtheorem{theorem}{Theorem}[section]
\newtheorem{proposition}[theorem]{Proposition}
\newtheorem{lemma}[theorem]{Lemma}
\newtheorem{corollary}[theorem]{Corollary}
\theoremstyle{definition}
\newtheorem{definition}[theorem]{Definition}
\theoremstyle{remark}
\newtheorem{remark}[theorem]{Remark}

\newcommand{\R}{\mathbb{R}}
\newcommand{\C}{\mathbb{C}}
\newcommand{\N}{\mathbb{N}}
\renewcommand{\Re}{\operatorname{Re}}
\newcommand{\GR}{\Gamma_{\R}}

\title{The information capacity of functional-equation windows}
\author{Samuel Lavery\thanks{Numerical instruments cited as data:
\texttt{tmp/adapter\_sym\_probe.py}, \texttt{tmp/adapter\_rigidity\_probe.py},
\texttt{tmp/adapter\_twist\_probe.py}, \texttt{tmp/adapter\_twist\_fast.py} in the
helix--Frobenius repository.  \texttt{sam@sk-labs.com}}}
\date{August 2026}

\begin{document}
\maketitle

\begin{abstract}
Fix an archimedean factor $G(s)=\prod_{i=1}^{m}\GR(s+\lambda_i)$ with
$0<\lambda_1\le\cdots\le\lambda_m$, its kernel
$K(y)=\frac{1}{2\pi i}\int_{(c)}G(s)\,y^{-s}\,ds$, and a compact window
$I=[a,b]\subset(0,\infty)$.  The \emph{windowed closure operator} sends a
coefficient vector $(c_n)_{n\le N}$ to the reflection defect
$x\mapsto\sum_n c_n\bigl(K(n/x)-\varepsilon\,x\,K(nx)\bigr)$ on $I$ --- the
linearization of the question ``which coefficient sequences satisfy this
functional equation on this window?''.  We prove three quantitative laws for
these operators and their character-twisted companions: exponential
singular-value decay with an explicit Bernstein rate (so a single functional
equation carries only $O(\log(1/\tau))$ constraint directions at relative
precision $\tau$, and the codimension of the windowed solution space is
correspondingly small); the resulting linear-in-$N$ deficiency and its
depth-stratification; and a \emph{burial law} for twists: twisting by a
primitive character of conductor $q$, which rescales the kernel argument by
$C=q^{(m_{\C}+\delta)/\dots}$, does not destroy the channel's information but
shifts its entire singular spectrum deeper by
$\lambda_1\log_{10}C$ decades, the minimal archimedean parameter acting as
the burial rate.  Numerical spectra for the symmetric-power stacks of the
discrete series of weight $12$ agree with all three laws, including the
predicted $7.9$-decade burial of the conductor-$3$ channel at the
$\mathrm{Sym}^5$ stack.  We record the interpretation for converse theory:
exact rigidity is a statement at infinite depth, per-channel information at
any finite depth is logarithmic, and conductor growth throttles twist
channels at a rate set by the archimedean type.
\end{abstract}

\section{Setup}\label{sec:setup}

Throughout, $\GR(s)=\pi^{-s/2}\Gamma(s/2)$ and
$G(s)=\prod_{i=1}^m \GR(s+\lambda_i)$ with real parameters
$0<\lambda_1\le\cdots\le\lambda_m$.  The associated kernel is
\[
K(y)\;=\;\frac{1}{2\pi i}\int_{(c)}G(s)\,y^{-s}\,ds\qquad(c>0),
\]
a Meijer $G$-function of type $G^{m,0}_{0,m}$ in the variable $\pi^m y^2$
\cite[Ch.~V]{Luke}; $K$ is holomorphic on the sector
$S_\phi=\{\,y:\ |\arg y|<\phi\,\}$ for any $\phi<m\pi/4$, positive on
$(0,\infty)$, of polynomial growth $y^{\lambda_1}$ as $y\to0^+$ and of
stretched-exponential decay
$|K(y)|\le C_K\exp(-c_m\,|y|^{2/m})$ in any proper subsector
\cite[Ch.~V, \S5.7--5.10]{Luke}.  Fix a window $I=[a,b]\subset(0,\infty)$, a
sign $\varepsilon\in\{\pm1\}$, and define, for $N\in\N\cup\{\infty\}$, the
\emph{windowed closure operator}
\[
A_N:\ \ell^2(\{1,\dots,N\})\longrightarrow L^2(I),\qquad
(A_Nc)(x)=\sum_{n\le N}c_n f_n(x),\quad
f_n(x)=K(n/x)-\varepsilon\,x\,K(nx).
\]
Since sectors are invariant under dilations and inversion, each $f_n$ extends
holomorphically to any fixed Bernstein ellipse
$E_\rho\subset S_\phi\cap\{\Re>0\}$ for $I$ (the image of $\{|w|=\rho\}$
under $w\mapsto\frac{b-a}{4}(w+w^{-1})+\frac{a+b}{2}$), and we set
$M_n=\sup_{E_\rho}|f_n|$; the kernel decay gives
$M_n\le C_\rho\exp(-c_\rho n^{2/m})$, so $\|M\|_2<\infty$ and $A_\infty$ is
well defined and compact.

\section{Window capacity}\label{sec:capacity}

\begin{theorem}[capacity of one functional equation]\label{thm:capacity}
Let $\rho>1$ be any Bernstein parameter with $E_\rho\subset
S_\phi\cap\{\Re>0\}$.  Then the singular values of $A_N$ (any $N\le\infty$)
satisfy, for all $k\ge0$,
\[
\sigma_{k+2}(A_N)\;\le\;\frac{2\sqrt{b-a}}{\rho-1}\,\|M\|_2\;\rho^{-k}.
\]
Consequently the relative $\tau$-rank obeys
\[
\#\{j:\ \sigma_j\ge\tau\,\sigma_1\}\;\le\;
2+\log_\rho\!\Bigl(\frac{2\sqrt{b-a}\,\|M\|_2}{(\rho-1)\,\tau\,\sigma_1}\Bigr)
\;=\;O\!\bigl(\log(1/\tau)\bigr).
\]
\end{theorem}

\begin{proof}
Let $c\in\ell^2$, $\|c\|_2\le1$, and $g=A_Nc$.  By the Cauchy--Schwarz
inequality $g$ extends holomorphically to $E_\rho$ with
$\sup_{E_\rho}|g|\le\sum_n|c_n|M_n\le\|M\|_2$.  Bernstein's theorem on
polynomial approximation of functions analytic in an ellipse
\cite[Thm.~8.2]{Trefethen} gives, for the space $P_k$ of polynomials of
degree $\le k$ (restricted to $I$, $\dim P_k=k+1$),
\[
\operatorname{dist}_{C(I)}(g,P_k)\;\le\;\frac{2\,\sup_{E_\rho}|g|}{\rho-1}\,
\rho^{-k}\;\le\;\frac{2\|M\|_2}{\rho-1}\,\rho^{-k},
\]
hence $\operatorname{dist}_{L^2(I)}(g,P_k)\le\sqrt{b-a}$ times the same
bound.  Thus the image of the unit ball of $A_N$ lies within
$\eta_k=\tfrac{2\sqrt{b-a}}{\rho-1}\|M\|_2\rho^{-k}$ of the
$(k{+}1)$-dimensional subspace $P_k\subset L^2(I)$, i.e.\ the
$(k{+}2)$-nd approximation number of $A_N$ is at most $\eta_k$; for compact
operators between Hilbert spaces approximation numbers coincide with
singular values \cite[Thm.~2.5.3]{Pinkus}, giving the claim.  The
$\tau$-rank bound is immediate.
\end{proof}

\begin{corollary}[deficiency, stratified by depth]\label{cor:deficiency}
For every relative floor $\tau$ and every $N$,
\[
\operatorname{nullity}_\tau(A_N)\;\ge\;N-2-
\log_\rho\!\Bigl(\frac{2\sqrt{b-a}\,\|M\|_2}{(\rho-1)\,\tau\,\sigma_1}\Bigr).
\]
The deficiency of a single windowed functional equation grows linearly in
the coefficient dimension, and the number of constraint directions grows
only logarithmically in the depth $1/\tau$: the ``hard band'' at shallow
depth is $O(1)$, and each further decade of depth releases at most
$1/\log_{10}\rho$ additional directions.
\end{corollary}

\begin{remark}
This is the theorem behind two measured phenomena: the fixed shallow
capacity of the bare functional equation across coefficient dimensions, and
the \emph{floor split} --- per-channel yields of a few directions at
$10^{-4}$ versus $\sim11$ at $10^{-20}$ --- which is the
$\log(1/\tau)$-slope made visible.
\end{remark}

\section{The burial law for twists}\label{sec:burial}

For a primitive character $\chi$ of conductor $q$ the twisted closure rows
carry the rescaled kernel: with $C$ the conductor scale of the twisted
completed function (for the rank-$d$ objects of \S\ref{sec:data},
$C=q^{d/2}$),
\[
f^{(q)}_n(x)\;=\;\chi(n)\bigl(K(n/(Cx))-\varepsilon_\chi\,x\,K(nx/C)\bigr),
\qquad A^{(q)}_N c=\sum_{n\le N}c_nf^{(q)}_n .
\]

\begin{theorem}[burial]\label{thm:burial}
Write the ascending exponent set of $K$ at the origin as
$\lambda_1<\lambda'<\cdots$ (the poles of $G(s)$, i.e.\ the numbers
$\lambda_i+2k$, listed without multiplicity), let
$r_0=\#\{i:\lambda_i=\lambda_1\}$, and put $\Delta=\lambda'-\lambda_1>0$.
Suppose $Nb\le C$ (every kernel argument on the window is $\le1$).  Then
\[
A^{(q)}_N\;=\;B_N\;+\;E_N,\qquad
\operatorname{rank}B_N\le 2r_0,\qquad
\|E_N\|\;\le\;C_1\,\Bigl(\frac{Nb}{C}\Bigr)^{\lambda_1+\Delta}\sqrt{N},
\]
where $B_N$ is built from the leading Frobenius term
$K(y)=\kappa_0y^{\lambda_1}(1+O(y^{\Delta}))$: its rows are
$\chi(n)\kappa_0 (n/C)^{\lambda_1}\bigl(x^{-\lambda_1}-\varepsilon_\chi
x^{1+\lambda_1}\bigr)$ up to the $r_0$-fold multiplicity structure.
Consequently the whole singular spectrum of the channel is suppressed by the
factor $(1/C)^{\lambda_1}$ relative to an unstarved channel: on a logarithmic
depth axis the channel's information is not destroyed but \emph{buried}
\[
\text{burial depth}\;\approx\;\lambda_1\,\log_{10}C\ \ \text{decades}.
\]
\end{theorem}

\begin{proof}
The Mellin--Barnes representation gives the convergent Frobenius expansion
of $K$ at the origin by shifting the contour left across the poles of
$G$ \cite[Ch.~V]{Luke}: $K(y)=\kappa_0y^{\lambda_1}+O(y^{\lambda_1+\Delta})$
uniformly on $0<y\le b$, with $r_0$ minimal-exponent terms when the minimal
parameter has multiplicity $r_0$ (log terms, absorbed in the same rank
count).  Substituting $y=n/(Cx)$ and $y=nx/C$ and collecting the leading
terms yields $B_N$, whose column space lies in the $\le2r_0$-dimensional
span of $x^{-\lambda_1},x^{1+\lambda_1}$ (times logs); the remainder rows
are bounded by $C_1(n/C)^{\lambda_1+\Delta}$ on $I$, and summing squares
over $n\le N$ gives the stated norm bound.  Weyl's inequality
$\sigma_{k+2r_0}(A^{(q)}_N)\le\|E_N\|$ then confines all but $2r_0$
directions below the remainder scale, while the surviving leading block
itself carries the prefactor $(1/C)^{\lambda_1}$; comparison with the
unrescaled channel gives the depth shift.
\end{proof}

\begin{remark}[the archimedean type sets the burial rate]
$\lambda_1$ is the minimal archimedean parameter of the stack.  For the
symmetric-power stacks of a weight-$12$ form, $\lambda_1=11/2$ at every odd
rank; at $\mathrm{Sym}^5$ (rank $6$, $C=q^3$) the predicted burial of the
conductor-$3$ channel is $\tfrac{11}{2}\log_{10}27\approx7.9$ decades.
\end{remark}

\section{Data}\label{sec:data}

All spectra are for the symmetric-power stacks of the discrete series of
weight $12$ on the window $I=[1,1.8]$, with exact symmetric-power
coefficients; floors are relative to $\sigma_1$.

\medskip
\noindent\emph{Capacity and floor split.}  Bare functional equation:
nullity $54/51/47$ of $60$ at rank $3$ and $34/31/28$ of $40$ at rank $6$
(floors $10^{-8}/10^{-14}/10^{-20}$); at matched shallow floors both ranks
carry $\sim3$ directions per channel, at $10^{-20}$ the rank-$3$ channels
carry $\sim11$ --- the $\log(1/\tau)$ slope of
Corollary~\ref{cor:deficiency}.

\medskip
\noindent\emph{Additivity at rank 3.}  Quadratic-twist channels
$q=3,4,5$ cut the rank-$3$ deficiency $47\to36\to24\to13$ (of $60$, floor
$10^{-20}$): $\sim11$ directions per channel, no saturation --- consistent
with the classical fact that $\mathrm{GL}(3)$'s converse theorem needs
$\mathrm{GL}(1)$ twists only.

\medskip
\noindent\emph{Starvation and burial at rank 6.}  At $N=120$ the shallow
yields of $q=3,4,5,7,8$ are $2,2,1,0,0$ --- monotone in the conductor scale
$C=q^3$, the shallow reading of Theorem~\ref{thm:burial}.  The deep
spectrum of the $q=3$ channel recovers $\sim10$ directions at $10^{-20}$
while yielding $\sim2$ at $10^{-4}$: the channel's information sits
$\approx8$ decades deep, matching the predicted
$\lambda_1\log_{10}C\approx7.9$.

\section{Interpretation for converse theory}\label{sec:interp}

Three consequences, stated as interpretation rather than theorem.  First,
\emph{exact} rigidity --- the converse-theorem regime --- is a statement at
infinite depth: any finite-precision reading of a functional-equation
system sees only logarithmically many constraints per channel, so the
empirical hardness of converse problems is structural, not numerical
accident.  Second, rigidifying $N$ coefficients at depth $\tau$ needs
total capacity $\ge N$, hence
$\#\text{channels}\gtrsim N/\bigl(O(\log(1/\tau))\bigr)$ \emph{before}
conductor effects; twist families are not a technical convenience of the
classical proofs but an information-theoretic necessity.  Third, the burial
law throttles high-rank twisting: the conductor scale $q^{d/2}$ grows with
the rank $d$, so at fixed depth the usable twist channels thin out
exponentially in the rank --- a quantitative shape for why the twist
requirements of converse theorems, and the difficulty of supplying them,
escalate with $d$.

\begin{remark}[geometric register]
The objects of this note are the one-dimensional chart readout of
carrier-native structure, and both laws state carrier facts.  The variable
$x$ is the carrier's scale flow ($\log x$ the height along the axis), and
$x\mapsto1/x$ is the reflection through the origin exchanging the two
conjugate ends of the double helix; the window is a height segment about the
self-dual point, and the closure defect is the end-to-end \emph{registration}
of the helix against the anti-helix on that segment.
Theorem~\ref{thm:capacity} says a finite height window of registration has
finite resolving power --- logarithmically many fiber channels per decade of
depth --- because the carrier's clock is smooth; the Bernstein ellipse is the
admissible tube around the axis, of width set by the rank's sector.
Theorem~\ref{thm:burial} is the phasor entry law made quantitative: the
exponent $\lambda_1$, a Mellin pole in the chart, is on the carrier the
growth law of a phasor at birth (each phasor enters at height zero and grows
continuously as $y^{\lambda_1}$), and a conductor-$C$ rail's phasors are
still in their entry regime at every height the base window sees --- its
information sits $\lambda_1\log_{10}C$ decades deep because information
lives where phasors have matured.
\end{remark}

\begin{remark}[scope and provenance]
The operators studied here impose the functional equation \emph{linearly};
no Euler-product or multiplicativity constraint is imposed, so nothing here
bounds the classical converse theorems themselves, whose hypotheses live at
exact depth with arithmetic structure.  The nearest existing body of
results known to us is the Selberg-class rigidity theory of Kaczorowski and
Perelli (classification at small degree from functional equation plus Euler
product); a systematic comparison is a stated obligation before any novelty
claim for the framework.  The numerical spectra are data, not proof; the
theorems above are unconditional statements about the operators.
\end{remark}

\begin{thebibliography}{9}

\bibitem{Luke} Y.~L. Luke, \emph{The Special Functions and Their
Approximations}, Vol.~I, Academic Press, 1969.

\bibitem{Trefethen} L.~N. Trefethen, \emph{Approximation Theory and
Approximation Practice}, SIAM, 2013.

\bibitem{Pinkus} A.~Pinkus, \emph{$n$-Widths in Approximation Theory},
Springer, 1985.

\bibitem{KP} J.~Kaczorowski, A.~Perelli, \emph{On the structure of the
Selberg class, I--VII}, Acta Math.\ \textbf{182} (1999) 207--241 and
sequels.

\end{thebibliography}

\end{document}
